%!TEX root = ../../main.tex

\chapter{Einleitung}
\label{ch:einleitung}


\section{Motivation und Problemstellung}
\label{subsec:motivation}

Die Verfügbarkeit, Qualität und zeitnahe Bereitstellung von Daten sind zu einem zentralen Erfolgsfaktor für datengetriebene Organisationen geworden. Im Zentrum der dafür notwendigen Dateninfrastruktur stehen seit Jahrzehnten \ac{ETL}-Prozesse, die Daten aus heterogenen Quellsystemen extrahieren, in eine einheitliche Struktur überführen und in Zielsysteme wie Data Warehouses oder Analyseplattformen laden \autocite{khanOverviewETLTechniques2024,walhaDataIntegrationTraditional2024}. Trotz ihrer langen Etablierung gelten klassische \ac{ETL}-Pipelines zunehmend als Engpass: Sie sind in der Regel regelbasiert, statisch modelliert und erfordern erheblichen manuellen Aufwand bei Entwurf, Anpassung und Wartung \autocite{walhaDataIntegrationTraditional2024,nagabhyruBridgingTraditionalETL2022}. Mit dem Zuwachs von Datenvolumen, der Heterogenität der Quellen und steigenden Anforderungen an Aktualität geraten herkömmliche \ac{ETL}-Architekturen an ihre Grenzen \autocite{walhaDataIntegrationTraditional2024,raviETLExtractTransform2025}.
Parallel dazu haben sich in den letzten Jahren \acp{LLM} als leistungsfähige Werkzeuge etabliert, die neben natürlicher Sprache auch strukturierte und semi-strukturierte Daten verarbeiten können \autocite{zhaoSurveyLargeLanguage2026}. Arbeiten zeigen, dass \acp{LLM} im Kontext der Datenverarbeitung Aufgaben wie Datenbereinigung, Schema-Mapping, Transformationslogik oder Anomalieerkennung übernehmen oder zumindest unterstützen können \autocite{hassaniRoleChatGPTData2023,mumuniAutomatedDataProcessing2025}. Damit eröffnet sich die Perspektive, klassische \ac{ETL}-Prozesse durch \ac{KI}-gestützte Komponenten zu ergänzen oder partiell zu ersetzen, um manuellen Aufwand zu reduzieren und die Anpassungsfähigkeit der Pipelines zu erhöhen. Der praktische Anreiz liegt dabei weniger in der technischen Machbarkeit an sich als vielmehr in der Aussicht auf Effizienz- und Kostenvorteile: Übernehmen Modelle Routineschritte der Transformationslogik ganz oder unterstützend, lässt sich der zeitintensive und damit kostenträchtige manuelle Entwicklungs- und Wartungsaufwand spürbar verringern, was den Einsatz gerade im wirtschaftlichen Kontext motiviert.
Gleichzeitig ist die praktische Eignung von \acp{LLM} für \ac{ETL}-Aufgaben bislang nur unzureichend systematisch untersucht. Während konzeptionelle Beiträge die potenziellen Vorteile beschreiben, fehlt es an empirischen Vergleichen, die unterschiedliche Modelle und Prompting-Strategien unter kontrollierten Bedingungen gegen klassische Verfahren stellen. Offen ist insbesondere, wie zuverlässig \acp{LLM} typische \ac{ETL}-Teilaufgaben wie Datenbereinigung, Datentransformation und Duplikaterkennung bearbeiten, wie stabil und reproduzierbar ihre Ergebnisse aufgrund der probabilistischen Modellarchitektur ausfallen, wie sich verschiedene Prompting-Ansätze auf die Ergebnisqualität auswirken und unter welchen Bedingungen ein Einsatz gegenüber etablierten regelbasierten Methoden sinnvoll ist. Genau an diesem Forschungsdefizit setzt die vorliegende Arbeit an. Sie betrachtet dabei den autonomen Einsatz der Modelle, also die Verarbeitung ohne menschliche Nachbearbeitung des erzeugten Ergebnisses. Damit bleibt eine weitere offene Frage bewusst außerhalb des Untersuchungsrahmens: wie sich ein solcher Einsatz zum heutigen De-facto-Standard der Softwareentwicklung verhält, bei dem \ac{KI} nicht autonom, sondern als Assistenz einer menschlichen Fachkraft wirkt. Die hier erhobenen Werte beschreiben folglich, was ohne Nachbearbeitung entsteht, und markieren damit die Untergrenze dessen, was im assistierten Einsatz zu erwarten wäre. Auf diese Abgrenzung wird in Kapitel~\ref{ch:fazit} zurückzukommen sein.
\section{Zielsetzung und Forschungsfragen}
\label{subsec:zielsetzung}

Ziel dieser Arbeit ist eine experimentelle Evaluierung, inwieweit sich aktuelle \acp{LLM} für ausgewählte \ac{ETL}-Aufgaben eignen, sei es als eigenständige Automatisierung oder als unterstützendes Werkzeug einer menschlichen Fachkraft. Im Zentrum steht der systematische Vergleich mehrerer aktueller \acp{LLM} und unterschiedlicher Prompting-Strategien, bewertet anhand von Genauigkeit, Vollständigkeit, Konsistenz und Effizienz bei der Lösung der definierten \ac{ETL}-Teilaufgaben und ausdrücklich nicht als allgemeiner Leistungsvergleich der Modelle. Grundlage ist ein kontrolliert erstellter synthetischer Datensatz, der typische Datenqualitätsprobleme abbildet. Die Ergebnisse werden mit einer klassisch implementierten \ac{ETL}-Pipeline verglichen, um Stärken, Schwächen und praktische Einsatzgrenzen \ac{KI}-gestützter Ansätze fundiert einordnen zu können. Auf dieser Grundlage sollen Handlungsempfehlungen für den Einsatz von \acp{LLM} in realen \ac{ETL}-Szenarien abgeleitet werden.
Aus dieser Zielsetzung ergeben sich zwei leitende Forschungsfragen:

\begin{description}
    \item[FF1:] Inwieweit eignen sich aktuelle \acp{LLM} für die Bearbeitung typischer \ac{ETL}-Teilaufgaben -- Datenbereinigung, Datentransformation und Duplikaterkennung -- im Vergleich zu einer klassischen, regelbasierten Pipeline, wie hängt diese Eignung neben dem Aufgabentyp von der Ausführungsform ab -- also davon, ob das Modell Verarbeitungscode erzeugt oder die Daten unmittelbar verarbeitet und wie die Aufgabe zugeschnitten ist --, und welches Vorgehen legen die Ergebnisse je Aufgabentyp nahe?
    \item[FF2:] Wie stark beeinflusst die Wahl der Prompting-Strategie die Ergebnisqualität der untersuchten \acp{LLM}, und welche Strategien sind für die betrachteten \ac{ETL}-Aufgaben zu empfehlen?
\end{description}

Da \acp{LLM} probabilistische Ausgaben erzeugen, wird die Reproduzierbarkeit ihrer Ergebnisse nicht als eigenständige Forschungsfrage, sondern als querschnittliche Untersuchungsdimension über beide Fragen hinweg behandelt: Jede Konfiguration wird mehrfach ausgeführt, sodass die Stabilität der Ergebnisse durchgängig mitbewertet werden kann.

\section{Methodisches Vorgehen}
\label{sec:methodik}

Die Arbeit folgt einem experimentellen Vorgehen: Nach einer literaturbasierten Fundierung wird ein kontrolliertes Versuchsdesign entwickelt, dessen Ergebnisse gegen eine regelbasiert implementierte \ac{ETL}-Pipeline als Maßstab gestellt werden. Grundlage aller Messungen ist ein synthetischer Datensatz mit gezielt eingebauten Qualitätsproblemen, dessen Soll-Lösung damit exakt bekannt ist. Gemessen wird auf mehreren Untersuchungsebenen, die nacheinander je einen Einflussfaktor freilegen: eine vollständig gekreuzte Faktormatrix aus Modell, Prompting-Strategie und Aufgabe; eine verkettete Pipeline aus neun aufeinander aufbauenden Schritten sowie ein sonst identischer isolierter Durchlauf, der eigene Schrittleistung von geerbtem Fehler trennt; ein Vergleich zweier Ausführungsformen, in denen das Modell entweder Verarbeitungscode erzeugt oder die Daten unmittelbar verarbeitet; und schließlich zwei Ebenen, die den Zuschnitt der Aufgabenstellung selbst zur Untersuchungsgröße machen. Da \acp{LLM} probabilistische Ausgaben erzeugen, wird jede Konfiguration mehrfach ausgeführt und die Streuung durchgängig mitberichtet. Die Auswertung verbindet die quantitativen Kennzahlen mit einer qualitativen Einordnung beobachteter Fehlermuster, aus der die abschließenden Handlungsempfehlungen abgeleitet werden.

\section{Aufbau der Arbeit}
\label{sec:aufbau}

Die vorliegende Arbeit gliedert sich in zehn Kapitel. Nach der vorliegenden Einleitung, in der Motivation, Zielsetzung, Forschungsfragen und methodisches Vorgehen dargelegt werden, behandelt Kapitel~\ref{ch:grundlagen} die theoretischen Grundlagen. Hier werden zentrale Konzepte zu \ac{ETL}-Prozessen, Datenqualität, \acp{LLM} sowie Prompt Engineering eingeführt, die als Bezugsrahmen für die weitere Untersuchung dienen.
Kapitel~\ref{ch:stand-der-forschung} widmet sich dem aktuellen Stand der Forschung. Es werden Arbeiten zum Einsatz von \ac{KI} in der Datenverarbeitung, zur Verarbeitung strukturierter Daten durch \acp{LLM} sowie zur Automatisierung von \ac{ETL}-Prozessen aufgearbeitet und gegeneinander abgegrenzt. Auf dieser Basis wird die für diese Arbeit relevante Forschungslücke präzisiert.
In Kapitel~\ref{ch:konzeption} wird die Konzeption der experimentellen Evaluierung beschrieben. Dies umfasst das Forschungsdesign, die Auswahl der untersuchten Modelle, die Definition der \ac{ETL}-Aufgaben, die eingesetzten Prompting-Strategien sowie die Evaluationskriterien einschließlich der manuellen Pipeline. Kapitel~\ref{ch:datensatz} dokumentiert anschließend die Erstellung des synthetischen Datensatzes, einschließlich Anforderungen, Datenmodell, eingebauter Qualitätsprobleme und der durchgeführten Validierung.
Die technische Umsetzung wird in Kapitel~\ref{ch:implementierung} dargestellt. Es beschreibt die Systemarchitektur, die eingesetzten Technologien sowie die konkrete Implementierung der klassischen \ac{ETL}-Pipeline und der \ac{LLM}-basierten Komponenten. Kapitel~\ref{ch:durchfuehrung} behandelt die Durchführung der Experimente und legt Versuchsablauf, erhobene Messdaten sowie die Bedingungen zur Sicherstellung der Reproduzierbarkeit offen.
Kapitel~\ref{ch:ergebnisse} präsentiert und analysiert die Ergebnisse. Neben einer aufgabenspezifischen Auswertung erfolgen ein Vergleich der Modelle, der Prompting-Strategien sowie eine Gegenüberstellung von \ac{KI}-basiertem und klassischem Ansatz. Die Ergebnisse werden anschließend diskutiert und in den Stand der Forschung eingeordnet. Aus den gewonnenen Erkenntnissen leitet Kapitel~\ref{ch:handlungsempfehlungen} Handlungsempfehlungen für die Praxis ab, insbesondere zur Eignung von \acp{LLM} in unterschiedlichen \ac{ETL}-Szenarien sowie zur Gestaltung hybrider Architekturen.
Kapitel~\ref{ch:fazit} fasst die zentralen Ergebnisse zusammen, beantwortet die in Abschnitt~\ref{subsec:zielsetzung} formulierten Forschungsfragen, reflektiert die Limitationen der Arbeit und gibt einen Ausblick auf weiterführende Forschungsfragen.
